% 第 11 题图丙：电压表表盘（0~3 V，指针示数 1.23 V）（TikZ 线稿）
% 分度值 0.1 V（每格读数需估读一位），指针位于 1.2 与 1.3 之间
\begin{tikzpicture}[font=\small, >=Stealth]
  % 表盘半圆弧（0 V 在左端 180°，3 V 在右端 0°）
  \draw[thick] (-2.2,0) arc[start angle=180, end angle=0, radius=2.2];
  \draw[thick] (-2.2,0) -- (2.2,0);
  % 小刻度（0~3 V 共 30 格，每格 0.1 V）
  \foreach \i in {1,...,29} {
    \pgfmathsetmacro\ang{180 - \i*6}
    \draw ({2.05*cos(\ang)},{2.05*sin(\ang)}) -- ({2.2*cos(\ang)},{2.2*sin(\ang)});
  }
  % 大刻度（0、1、2、3 V）
  \foreach \i in {0,10,20,30} {
    \pgfmathsetmacro\ang{180 - \i*6}
    \draw[thick] ({1.95*cos(\ang)},{1.95*sin(\ang)}) -- ({2.2*cos(\ang)},{2.2*sin(\ang)});
  }
  % 数字刻度 0、1、2、3
  \node at (-2.42,0.12) {0};
  \node at (-0.78,1.34) {1};
  \node at (0.78,1.34) {2};
  \node at (2.42,0.12) {3};
  % 指针（示数 1.23 V：对应角度 180-6×12.3 = 106.2°）
  \draw[very thick, ->] (0,0) -- ({2.15*cos(106.2)},{2.15*sin(106.2)});
  \fill (0,0) circle (2pt);
  % 表盘标注
  \node[below] at (0,-0.5) {量程 0~3 V，示数 1.23 V};
\end{tikzpicture}
